\section{Kinematic Corrections in Fake Estimation}
\label{sec:kinematic-corrections}
The MC samples of a reference channel are used to derive this correction. %
Notice that a `bin-by-bin correction' approach cannot be applied since a given value of CB $p_{\text{T}}$ corresponds to multiple values of ID $p_{\text{T}}$. %
Thus, bin-by-bin corrections derived using one sample will not be applicable to another sample. %
Rather, we adopt the following approach:
\begin{itemize}
    \item For each event in the reference channel, calculate the response of $p_{\text{T}}$ of the two leptons:
    \begin{align}
        \mathcal{R}_{p^+_{\text{T}}} &= \frac{p^+_{\text{T}}({\text{CB}}) - p^+_{\text{T}}({\text{ID}})}{p^+_{\text{T}}({\text{ID}})}\\
        \mathcal{R}_{p^-_{\text{T}}} &= \frac{p^-_{\text{T}}({\text{CB}}) - p^-_{\text{T}}({\text{ID}})}{p^-_{\text{T}}({\text{ID}})}
    \end{align}
    \item Choose a binning in the $3{\text{D}}$ space of the $p_{\text{T}}$ of the two leptons and the angular distance between the two leptons, $\Delta R$ -- where all three variables are calculated using the {\text{ID}} tracks. For each bin $B$ (say indexed by $(p,q,r)$) in this $3{\text{D}}$ space, compile an unordered list of $2{\text{D}}$ responses $L_R(p,q,r) = \big[\big(\mathcal{R}^{i}_{p^+_{\text{T}}}, \mathcal{R}^{i}_{p^-_{\text{T}}}\big)\ \big\vert\  i\in B_{pqr}\big]$ -- where $i$ indexes the events in the reference channel. Correspondingly, for each bin, compile an unbinned list of weights $L_W(p,q,r)=\big[W^i=\frac{w^i}{\sum_j w^j}\ \big\vert i,j\in B_{pqr}\big]$, where $w^i$ is the weight of the event $i$ in the reference channel. 
    \item For each event $j$ in the target sample, i.e., the sample for which we want to apply the kinematic corrections, calculate the unbinned list of corrected $p_{\text{T}}$ values using the unbinned list of responses compiled in the previous step -- depending on the bin of $(p^+_{\text{T}}, p^-_{\text{T}}, \Delta R)$ that the event lies in. To be explicit, the list of corrected $p_{\text{T}}$ of the two leptons is calculated as:
    \begin{align}
        p^+_{\text{T}}({\text{Corrected}})_{ij} &= p^+_{\text{T}}({\text{ID}})_j \times (1 + \mathcal{R}^{i}_{p^+_{\text{T}}}(p(p^+_{\text{T}}({\text{ID}})_j),q(p^-_{\text{T}}({\text{ID}})_j),r(\Delta R({\text{ID}})_j)))\\
        p^-_{\text{T}}({\text{Corrected}})_{ij} &= p^-_{\text{T}}({\text{ID}})_j \times (1 + \mathcal{R}^{i}_{p^-_{\text{T}}}(p(p^+_{\text{T}}({\text{ID}})_j),q(p^-_{\text{T}}({\text{ID}})_j),r(\Delta R({\text{ID}})_j)))
    \end{align}
    where we use the notation $p(p^+_{\text{T}}({\text{ID}})_j)$ to denote the bin of $p^+_{\text{T}}({\text{ID}})_j$ in the $3{\text{D}}$ space of $(p^+_{\text{T}}, p^-_{\text{T}}, \Delta R)$, and the notation of $\mathcal{R}_{p^+_{\text{T}}}^i(p,q,r)$ to denote the $i^{\text{th}}$ $\mathcal{R}_{p^+_{\text{T}}}$ response in the list $L_R(p,q,r)$. % 
    \item Following a similar notation, the weight associated with the $i^{\text{th}}$ correction for the $j^{\text{th}}$ event in the target sample is given by:
    \begin{align}
        W_{ij} &= W^i(p(p^+_{\text{T}}({\text{ID}})_j),q(p^-_{\text{T}}({\text{ID}})_j),r(\Delta R({\text{ID}})_j))
    \end{align}
    \item Finally, the corrected $p_{\text{T}}$s of the two leptons of the target sample are used to calculate the kinematic variables of interest, e.g., the mass of the $B$-candidate, the $p_{\text{T}}$ of the $B$-candidate, etc. 
\end{itemize}

\subsection{Kinematic correction results}
\label{sec:kinematic-correction-results}

The corrections were calculated for 9 $\Delta R $ bins. Due to the non-uniform accumulation of data in the different $\Delta R$ bins, there was used variable bin sizes to include considerable amount of data in regions with low statistics and higher granularity in regions with sufficient amount of data. 

\begin{figure}[H] % <-- Use [H] here
  \centering
  \begin{minipage}[t]{0.38\textwidth}
    \centering
    \foreach \pt in {1,3,5,7,9}{
      \begin{subfigure}%{0.95\textwidth}
        \includegraphics[width=0.9\textwidth]{figures/KinematicCorrections/pT1_pT2_DR\pt}
        %\caption*{$p_T=\pt$}
      \end{subfigure}
    }
  \end{minipage}
  \hfill
  \begin{minipage}[t]{0.38\textwidth}
    \centering
    \foreach \pt in {2,4,6,8}{
      \begin{subfigure}%{0.95\textwidth}
        \includegraphics[width=0.9\textwidth]{figures/KinematicCorrections/pT1_pT2_DR\pt}
        %\caption*{$p_T=\pt$, $\eta=2$}
      \end{subfigure}
    }
  \end{minipage}
  \caption{Correction factors for different $\Delta R$ bins. }
\end{figure}